% Interfaz computable CG-008; formalización nueva de un lector residual.
% Productor: los caracteres e intervalos regionales construidos previamente.
% No identifica la componente numérica con el estado enriquecido completo.

\section{Inicialización residual computable y compatibilidad de las emisiones regionales}
\label{sec:residuo-regional-computable}

La realización numérica de una sección regional admite una descripción
ejecutable por intervalos racionales. Esta descripción permite inicializar
el residuo, actualizarlo y recuperar sus bloques sin introducir una cadena
decimal previamente fijada. Su origen es el productor regional del
\cref{x_reciprocidad:sec:generacion}: primero quedan determinados el carácter,
su orientación y su normalización; después se construyen sus aproximaciones
racionales. El presente desarrollo formaliza la componente residual de esa
lectura. La trayectoria, la hoja, las incidencias y la memoria del estado
enriquecido conservan simultáneamente sus operaciones propias.

El orden de las operaciones es, por tanto,
\[
 \begin{gathered}
 \text{estado regional generado}
 \longrightarrow \text{intervalos racionales del carácter}
 \\
 \longrightarrow \text{nombre racional del residuo}
 \longrightarrow \text{emisiones posicionales compatibles}.
 \end{gathered}
\]
La comparación con la realización arquimediana aparece en la demostración
de corrección. El algoritmo recibe intervalos; ninguna de sus instrucciones
recibe un número real convencional ni una lista de cifras objetivo.

\subsection{Nombre racional producido por un lector regional}
\label{subsec:residual-nombre-racional}

Se fija una región y se escribe \(j\mapsto[\ell_j,u_j]\) para su productor
racional. Las condiciones de la interfaz son
\begin{equation}
 \ell_j\leq v\leq u_j,\qquad
 \lim_{j\to\infty}(u_j-\ell_j)=0,
 \qquad \ell_j,u_j\in\mathbb Q.
 \label{eq:residual-contrato-productor}
\end{equation}
Aquí \(v\) designa la realización exacta del carácter para formular la
corrección; el procedimiento que calcula \(\ell_j,u_j\) procede del lector
regional. Los intervalos de clausura, propagación y autoescala ya
construidos satisfacen estas condiciones. En la implementación, su interfaz
es una función que, dado \(j\), devuelve dos fracciones exactas.

La anidación se obtiene, si hace falta, mediante la operación finita
\begin{equation}
 L_j=\max_{0\leq i\leq j}\ell_i,
 \qquad U_j=\min_{0\leq i\leq j}u_i,
 \qquad I_j=[L_j,U_j].
 \label{eq:residual-anidacion}
\end{equation}

\begin{proposition}[Anidación y conservación de la realización]
\label{prop:residual-anidacion}
Los intervalos \(I_j\) son racionales, no vacíos y encajados. Contienen
\(v\), su anchura tiende a cero y \(\bigcap_j I_j=\{v\}\).
\end{proposition}
\begin{proof}
Cada cota inferior es menor o igual que \(v\), y cada cota superior es
mayor o igual que \(v\); la misma propiedad vale para el máximo y el mínimo
finitos. Al aumentar \(j\), el máximo crece y el mínimo decrece.
Además, \(0\leq U_j-L_j\leq u_j-\ell_j\). Si dos puntos pertenecieran a
la intersección, su distancia estaría acotada por cada anchura y sería cero.
La pertenencia de \(v\) acredita la no vacuidad.
\end{proof}

\subsection{Prefijo inicial y representación del residuo}
\label{subsec:residual-inicializacion}

Sean \(b\geq2\) una base entera y \(n_0\geq0\) una profundidad.
Se supone que la realización evita las fronteras correspondientes:
\begin{equation}
 b^n v\notin\mathbb Z\qquad(n\geq0).
 \label{eq:residual-no-frontera}
\end{equation}
La irracionalidad de los tres valores regionales implica esta propiedad;
su demostración o su empleo como hipótesis pertenece al lector fijado.
Para una profundidad concreta basta la evitación en las escalas que se
utilizan. Las realizaciones con expansión terminal requieren, adicionalmente,
un decisor exacto de frontera y una convención de representación.

Definimos el predicado racional
\[
 \mathsf{Sep}(j,n):\quad
 \lfloor b^n L_j\rfloor=\lceil b^n U_j\rceil-1.
\]
Es decidible mediante operaciones enteras sobre numeradores y
denominadores. Sean
\begin{equation}
 j_0=\min\{j:\mathsf{Sep}(j,n_0)\},
 \qquad N_0=\lfloor b^{n_0}L_{j_0}\rfloor.
 \label{eq:residual-prefijo-inicial}
\end{equation}
El entero \(N_0\) se calcula desde el productor. Su uso como antecedente
de la representación residual no exige introducirlo como dato literal.

\begin{theorem}[Inicialización computable del residuo regional]
\label{thm:residual-inicializacion}
El índice \(j_0\) existe. Para \(j\geq j_0\), los intervalos
\begin{equation}
 E_{0,j}=
 \bigl[b^{n_0}L_j-N_0,\ b^{n_0}U_j-N_0\bigr]
 \label{eq:residual-nombre-inicial}
\end{equation}
son no vacíos, encajados y contenidos en \([0,1]\). Su anchura tiende a
cero y su intersección consiste en un único elemento
\(\varepsilon_0\in(0,1)\). En la realización arquimediana se cumple
\begin{equation}
 N_0=\lfloor b^{n_0}v\rfloor,
 \qquad b^{n_0}v=N_0+\varepsilon_0.
 \label{eq:residual-identificacion-inicial}
\end{equation}
La sucesión racional \((E_{0,j})_{j\geq j_0}\) constituye una
representación computable del residuo.
\end{theorem}
\begin{proof}
El punto \(b^{n_0}v\) está en el interior de una celda de extremos
enteros. Su distancia a esos extremos es positiva. Cuando la anchura de
\(b^{n_0}I_j\) es menor que esa distancia, ambos extremos pertenecen a la
misma celda y se cumple \(\mathsf{Sep}(j,n_0)\). Existe por ello un
primer índice de separación. En cualquier índice separado, el entero
\(\lfloor b^{n_0}v\rfloor\) se encuentra entre los dos enteros iguales
del predicado y coincide con ellos. La igualdad inicial sitúa los extremos
de \(b^{n_0}I_{j_0}\) entre \(N_0\) y \(N_0+1\). La anidación
preserva esa inclusión. La transformación afín de
\eqref{eq:residual-nombre-inicial} conserva anidación y no vacuidad y
multiplica las anchuras por \(b^{n_0}\). Su intersección es el punto de
\eqref{eq:residual-identificacion-inicial}; la evitación de frontera lo
sitúa en \((0,1)\). Cada extremo se calcula mediante fracciones exactas,
y la búsqueda termina por el argumento anterior.
\end{proof}

La identidad \eqref{eq:residual-identificacion-inicial} es una conclusión
de corrección. La representación ejecutable es la sucesión
\eqref{eq:residual-nombre-inicial}, construida a partir de intervalos
regionales. Esta distinción separa la producción de un nombre racional del
reconocimiento de su valor.

\subsection{Actualización residual y emisión mediante aritmética exacta}
\label{subsec:residual-actualizacion}

Supongamos construido el estado de lectura \((n_0+k,N_k,j_k)\), junto
con el nombre
\begin{equation}
 E_{k,j}=
 \bigl[b^{n_0+k}L_j-N_k,\ b^{n_0+k}U_j-N_k\bigr],
 \qquad j\geq j_k.
 \label{eq:residual-nombre-k}
\end{equation}
Se aumenta \(j\) desde \(j_k\) hasta que
\[
 \left\lfloor b\inf E_{k,j}\right\rfloor
 =\left\lceil b\sup E_{k,j}\right\rceil-1.
\]
El primer índice que satisface esta igualdad se denota \(j_{k+1}\), y el
entero común se denota \(d_k\). La actualización es
\begin{equation}
 N_{k+1}=bN_k+d_k,
 \qquad E_{k+1,j}=bE_{k,j}-d_k
 \quad(j\geq j_{k+1}).
 \label{eq:residual-actualizacion-exacta}
\end{equation}
La multiplicación de un intervalo por \(b>0\) y la resta de \(d_k\)
actúan sobre sus dos extremos. Se conservan el índice de lectura alcanzado,
la profundidad y el bloque emitido.

\Needspace{20\baselineskip}
\begin{theorem}[Terminación a profundidad arbitraria y conmutación posicional]
\label{thm:residual-compatibilidad-universal}
Bajo \eqref{eq:residual-contrato-productor} y
\eqref{eq:residual-no-frontera}, cada actualización termina y, para todo
\(k\geq0\), se cumplen
\begin{align}
 N_k&=\lfloor b^{n_0+k}v\rfloor,
 &0&\leq d_k<b,\label{eq:residual-prefijos-k}\\
 \varepsilon_{k+1}&=b\varepsilon_k-d_k,
 &d_k&=\lfloor b\varepsilon_k\rfloor,\label{eq:residual-recurrencia}\\
 b^{n_0+k}v&=N_k+\varepsilon_k,
 &0&<\varepsilon_k<1.\label{eq:residual-invariante}
\end{align}
Si \(P(b,n)\) es la emisión por separación racional del
\cref{thm:df-emision-correcta}, entonces
\begin{equation}
 N_k=P(b,n_0+k),\qquad
 d_k=P(b,n_0+k+1)\bmod b.
 \label{eq:residual-igualdad-publicacion}
\end{equation}
Los índices de parada de ambas implementaciones pueden ser distintos;
sus enteros emitidos coinciden exactamente.
\end{theorem}
\begin{proof}
La inicialización proporciona las afirmaciones para \(k=0\). Supuestas
en el nivel \(k\), la multiplicación del invariante por \(b\) da
\[
 b^{n_0+k+1}v=bN_k+b\varepsilon_k.
\]
La evitación de frontera implica \(b\varepsilon_k\notin\mathbb Z\).
Los intervalos \(bE_{k,j}\) son encajados, contienen ese punto y tienen
anchura tendente a cero. El argumento de separación del teorema anterior
garantiza la terminación y fija
\(d_k=\lfloor b\varepsilon_k\rfloor\). Como
\(0<\varepsilon_k<1\), el bloque pertenece a
\(\{0,\ldots,b-1\}\). Restar \(d_k\) transforma los intervalos en
un nombre encajado del residuo siguiente y conserva su pertenencia a
\((0,1)\). La parte entera de la igualdad precedente es
\(bN_k+d_k\), lo que prueba el nuevo invariante. La inducción vale para
cada profundidad finita, sin imponer una profundidad máxima.
Finalmente, el teorema de emisión racional identifica también
\(P(b,n)\) con \(\lfloor b^n v\rfloor\). La igualdad de los
prefijos y la fórmula del residuo módulo \(b\) demuestran
\eqref{eq:residual-igualdad-publicacion}.
\end{proof}

\begin{corollary}[Truncamiento, recuperación y cotas residuales]
\label{cor:residual-recuperacion}
Para \(k,h\geq0\),
\[
 \left\lfloor\frac{N_{k+h}}{b^h}\right\rfloor=N_k,
 \qquad N_k=\frac{N_{k+1}-d_k}{b},
 \qquad \varepsilon_k=\frac{\varepsilon_{k+1}+d_k}{b}.
\]
La anchura del nombre residual al índice \(j\) es
\(b^{n_0+k}(U_j-L_j)\). En consecuencia, cualquier precisión racional
positiva para el residuo se alcanza tras un número finito de consultas al
productor. La profundidad y los bloques archivados determinan la
recuperación en la componente de lectura.
\end{corollary}
\begin{proof}
La división del prefijo extenso recupera el prefijo anterior por
\eqref{eq:residual-prefijos-k}. Las dos fórmulas de recuperación se
obtienen despejando en \eqref{eq:residual-actualizacion-exacta} y
\eqref{eq:residual-recurrencia}. La fórmula de anchura es inmediata de
\eqref{eq:residual-nombre-k}; su convergencia a cero proporciona la
terminación de la búsqueda de precisión.
\end{proof}

\subsection{Correspondencia de profundidades ternarias y decimales}
\label{subsec:residual-bases}

En la normalización fraccionaria se escribe \(v=m+r\), donde
\(m=P(b,0)\) se obtiene del mismo productor. El prefijo fraccionario es
\(C_n=P(b,n)-mb^n\). Por cálculo entero,
\[
 b^n v-P(b,n)=b^n r-C_n.
\]
Esta igualdad reconcilia el prefijo que incluye la parte entera con el que
se inicia en cero, conservando idéntico residuo. La resta de \(m\)
es una transformación posterior del nombre regional.

La agrupación trítica usa \(b=729=3^6\). Seis bloques corresponden a
\(36\) trits; la especialización \(n_0=6\) proporciona un nombre
racional del residuo de esa profundidad. La agrupación decimal usa
\(b=1000=10^3\); seis bloques corresponden a \(18\) cifras decimales.
Son coordenadas de lectura diferentes del mismo carácter. Su compatibilidad
se expresa por los cilindros que contienen esa realización y por las
identidades de cambio de escala; no exige igualar sus residuos numéricos.

El antecedente residual de
\cref{sec:c27-generacion-coinductiva-profundidad-arbitraria} distingue el
estado interno de su publicación. La presente interfaz aporta una
representación racional ejecutable de la componente numérica, con la misma
recurrencia residual una vez fijada su lectura. Para identificarla con otra
realización residual de esa sección basta acreditar que ambas leen el mismo
carácter, a la misma profundidad y con igual normalización: la identidad
\(b^n v=N_n+\varepsilon_n\) y la unicidad del prefijo fuerzan entonces
la igualdad de los residuos. Esta comparación es una igualdad de lectores.
La evolución de posiciones, hojas, incidencia y memoria permanece regida
por el TPK enriquecido y por sus respectivos teoremas de transporte.

\subsection{Contrato de la implementación exacta}
\label{subsec:residual-implementacion}

La implementación de esta interfaz emplea exclusivamente fracciones
racionales, potencias enteras, máximo, mínimo, división con suelo y división
con techo. El productor de intervalos es un argumento explícito del
programa. Una consulta no devuelve \(v\): devuelve sus dos cotas generadas.
El estado de ejecución conserva base, profundidad, prefijo, índice alcanzado
y registro ordenado de emisiones. Para cada paso, el registro guarda el
nombre residual anterior en el índice decisivo, el bloque y el nombre
transformado; las dos recuperaciones del corolario se comprueban mediante
igualdades racionales.

Las pruebas universales anteriores justifican la terminación y la
compatibilidad bajo el contrato del productor. Los ensayos finitos verifican
la implementación de esas operaciones y sus identidades; la generalidad
del resultado reside en la inducción y en la separación racional de celdas.
Esta formalización reúne una interfaz computable de lectura y recuperación
sin sustituir la selección regional ni atribuir a la componente numérica
las operaciones adicionales del estado enriquecido.
